\section{Introduction}
\label{sec:intro}

Deep generative models learn a distribution $p_\theta(\x)$ from data and can draw new samples from it. Two
families laid the foundations of modern image generation:
\begin{itemize}
  \item \textbf{Variational Autoencoders} (VAEs)~\citep{kingma2014vae} fit a latent-variable model by maximising
        a variational lower bound with an amortised inference network.
  \item \textbf{Denoising Diffusion Probabilistic Models} (DDPMs)~\citep{sohl2015deep,ho2020ddpm} learn to invert
        a fixed, gradual noising process.
\end{itemize}
Both are trained by maximising an evidence lower bound (ELBO), yet they behave very differently in practice.

This report implements both models from first principles, trains them on CIFAR-10~\citep{krizhevsky2009cifar},
and compares them quantitatively and qualitatively, with the goal of explaining the \emph{main difference}
between the two approaches.

\paragraph{Contributions.}
\begin{enumerate}
  \item \textbf{Faithful implementations from scratch} (\cref{sec:arch,sec:training}). The DDPM reproduces the
        paper's CIFAR-10 configuration exactly (35.7M-parameter U-Net); only the training length is shortened.
        The VAE follows the original method exactly and is instantiated with two networks: the paper's MLP and a
        convolutional network built from the U-Net's blocks.
  \item \textbf{A shared, validated evaluation protocol} (\cref{sec:eval}). It covers FID, IS, KID,
        precision/recall, likelihood, the speed--quality trade-off, classifier-based thematic consistency,
        memorisation and a blind rating. Before any model was evaluated, the metrics were checked to reproduce
        published reference values on real data.
  \item \textbf{A mechanistic comparison} (\cref{sec:results,sec:qualitative,sec:discussion}). It links each
        empirical difference (blur, diversity, the ``prior hole'', likelihood, sampling cost) to the structure of
        the two models.
\end{enumerate}

\paragraph{Dataset.} We use CIFAR-10~\citep{krizhevsky2009cifar}: 50,000 training and 10,000 test images of
$32\times32$ RGB pixels in 10 classes. It is the DDPM paper's main benchmark, so published numbers exist to check our
implementation against. Its classes make thematic consistency measurable, and the Inception Score is only
meaningful on ImageNet-like natural images such as these. Every training run and every speed measurement used one
NVIDIA A100-SXM4-80GB (Google Colab), so compute and timings are directly comparable.

\paragraph{Models compared.} Throughout the report we compare three models:
\begin{itemize}
  \item \textbf{VAE-MLP}: the VAE with the network of \citet{kingma2014vae}.
  \item \textbf{VAE-Conv}: the same VAE method with a convolutional network of comparable size to the DDPM. This
        is the fair counterpart to the DDPM.
  \item \textbf{DDPM}: the model of \citet{ho2020ddpm}.
\end{itemize}

\paragraph{Code and data availability.} All code, configurations, training logs, raw evaluation outputs and
figures are public at \repo. Every number in this report is produced by the scripts in that repository.
